\documentclass[12pt,a4paper]{article}
\usepackage[UTF8]{ctex}
\usepackage{amsmath,amssymb,amsthm}
\usepackage{geometry}

\geometry{margin=2.5cm}

\title{阻抗控制的两种实现方式：Twist/Wrench vs 解耦}
\author{第11章补充说明}
\date{}

\begin{document}

\maketitle

\section{问题提出}

在阻抗控制中，行为(11.64)有两种实现方式：

\begin{enumerate}
\item \textbf{使用twists和wrenches}（不解耦）：
\begin{itemize}
\item 将$f_{\text{ext}}$替换为wrench $F_{\text{ext}}$
\item 将$\dot{x}$替换为twist $V$
\item 将$\ddot{x}$替换为$\dot{V}$
\item 将$x$替换为指数坐标$S\theta$
\end{itemize}

\item \textbf{解耦方式}：
\begin{itemize}
\item 将线性和旋转行为分开处理
\end{itemize}
\end{enumerate}

\textbf{问题}：
\begin{enumerate}
\item 前者（使用twists和wrenches）是否解耦？
\item 两种方法有什么区别？
\item 各自的优缺点是什么？
\end{enumerate}

\section{原始方程回顾}

\subsection{方程(11.64)}

\textbf{原始阻抗控制方程}：
\begin{equation}
M\ddot{x} + B\dot{x} + Kx = f_{\text{ext}} \tag{11.64}
\end{equation}

其中：
\begin{itemize}
\item $x \in \mathbb{R}^n$：任务空间配置（最小坐标集）
\item 例如：$x \in \mathbb{R}^3$（只有位置，没有旋转）
\item $M, B, K$：虚拟质量、阻尼、刚度矩阵
\item $f_{\text{ext}}$：外力
\end{itemize}

\textbf{限制}：
\begin{itemize}
\item 这个方程只适用于\textbf{最小坐标集}
\item 例如：只有位置$x \in \mathbb{R}^3$，没有旋转
\item 不能直接表示SE(3)中的完整6维位姿
\end{itemize}

\section{方法1：使用Twists和Wrenches（耦合方式）}

\subsection{替换规则}

\textbf{替换}：
\begin{align}
f_{\text{ext}} &\to F_{\text{ext}} \quad \text{（wrench，6维）} \\
\dot{x} &\to V \quad \text{（twist，6维）} \\
\ddot{x} &\to \dot{V} \quad \text{（twist的导数，6维）} \\
x &\to S\theta \quad \text{（指数坐标，6维）}
\end{align}

\textbf{注意}：
\begin{itemize}
\item $F_{\text{ext}}$可以是体坐标系wrench $F_b$或空间坐标系wrench $F_s$
\item $V$可以是体坐标系twist $V_b$或空间坐标系twist $V_s$
\item $S\theta$是SE(3)的指数坐标表示
\end{itemize}

\subsection{替换后的方程}

\textbf{使用twists和wrenches的阻抗控制方程}：
\begin{equation}
M\dot{V} + BV + K(S\theta) = F_{\text{ext}} \tag{11.64-twist}
\end{equation}

其中：
\begin{itemize}
\item $M \in \mathbb{R}^{6 \times 6}$：虚拟质量矩阵（6维）
\item $B \in \mathbb{R}^{6 \times 6}$：虚拟阻尼矩阵（6维）
\item $K(S\theta)$：虚拟刚度项（依赖于指数坐标$S\theta$）
\item $F_{\text{ext}} \in \mathbb{R}^6$：外力wrench（6维）
\item $V \in \mathbb{R}^6$：twist（6维）
\item $\dot{V} \in \mathbb{R}^6$：twist的导数（6维）
\end{itemize}

\subsection{是否解耦？}

\textbf{答案}：\textcolor{red}{不解耦}！

\textbf{原因}：
\begin{itemize}
\item 矩阵$M, B, K$是$6 \times 6$的，通常是\textbf{非对角}的
\item 旋转和平移之间存在\textbf{耦合}
\item 例如：旋转运动可能影响平移，平移运动可能影响旋转
\end{itemize}

\textbf{耦合的例子}：
\begin{itemize}
\item 如果$M$不是对角矩阵，旋转加速度$\dot{\omega}$会影响平移加速度$\dot{v}$
\item 如果$B$不是对角矩阵，旋转速度$\omega$会影响平移速度$v$
\item 如果$K$不是对角矩阵，旋转位置$S\theta_{\text{rot}}$会影响平移位置$S\theta_{\text{trans}}$
\end{itemize}

\subsection{矩阵结构}

\textbf{一般形式}（耦合）：
\begin{equation}
M = \begin{bmatrix}
M_{\omega\omega} & M_{\omega v} \\
M_{v\omega} & M_{vv}
\end{bmatrix} \in \mathbb{R}^{6 \times 6}
\end{equation}

其中：
\begin{itemize}
\item $M_{\omega\omega} \in \mathbb{R}^{3 \times 3}$：旋转-旋转耦合
\item $M_{\omega v} \in \mathbb{R}^{3 \times 3}$：旋转-平移耦合
\item $M_{v\omega} \in \mathbb{R}^{3 \times 3}$：平移-旋转耦合
\item $M_{vv} \in \mathbb{R}^{3 \times 3}$：平移-平移耦合
\end{itemize}

\textbf{如果$M_{\omega v} \neq 0$或$M_{v\omega} \neq 0$}：
\begin{itemize}
\item 旋转和平移是耦合的
\item 旋转运动会影响平移，平移运动会影响旋转
\end{itemize}

\section{方法2：解耦方式}

\subsection{基本思想}

\textbf{解耦}：将线性和旋转行为分开处理。

\textbf{方法}：
\begin{itemize}
\item 将6维问题分解为两个独立的3维问题
\item 线性部分（平移）：3维
\item 旋转部分：3维
\item 两者独立处理，互不影响
\end{itemize}

\subsection{解耦后的方程}

\textbf{线性部分（平移）}：
\begin{equation}
M_{\text{trans}}\ddot{p} + B_{\text{trans}}\dot{p} + K_{\text{trans}}p = f_{\text{ext,trans}} \tag{11.64-trans}
\end{equation}

其中：
\begin{itemize}
\item $p \in \mathbb{R}^3$：位置
\item $\dot{p} \in \mathbb{R}^3$：线速度
\item $\ddot{p} \in \mathbb{R}^3$：线加速度
\item $M_{\text{trans}}, B_{\text{trans}}, K_{\text{trans}} \in \mathbb{R}^{3 \times 3}$：平移的虚拟质量、阻尼、刚度矩阵
\item $f_{\text{ext,trans}} \in \mathbb{R}^3$：外力（平移部分）
\end{itemize}

\textbf{旋转部分}：
\begin{equation}
M_{\text{rot}}\dot{\omega} + B_{\text{rot}}\omega + K_{\text{rot}}\theta_{\text{rot}} = m_{\text{ext}} \tag{11.64-rot}
\end{equation}

其中：
\begin{itemize}
\item $\theta_{\text{rot}} \in \mathbb{R}^3$：旋转角度（或指数坐标的旋转部分）
\item $\omega \in \mathbb{R}^3$：角速度
\item $\dot{\omega} \in \mathbb{R}^3$：角加速度
\item $M_{\text{rot}}, B_{\text{rot}}, K_{\text{rot}} \in \mathbb{R}^{3 \times 3}$：旋转的虚拟质量、阻尼、刚度矩阵
\item $m_{\text{ext}} \in \mathbb{R}^3$：外力矩（旋转部分）
\end{itemize}

\subsection{解耦的矩阵结构}

\textbf{解耦的虚拟质量矩阵}：
\begin{equation}
M = \begin{bmatrix}
M_{\text{rot}} & 0 \\
0 & M_{\text{trans}}
\end{bmatrix} \in \mathbb{R}^{6 \times 6}
\end{equation}

其中：
\begin{itemize}
\item $M_{\text{rot}} \in \mathbb{R}^{3 \times 3}$：旋转质量矩阵（对角或非对角，但只影响旋转）
\item $M_{\text{trans}} \in \mathbb{R}^{3 \times 3}$：平移质量矩阵（对角或非对角，但只影响平移）
\item 非对角块为零：$M_{\omega v} = 0$，$M_{v\omega} = 0$
\end{itemize}

\textbf{解耦的虚拟阻尼矩阵}：
\begin{equation}
B = \begin{bmatrix}
B_{\text{rot}} & 0 \\
0 & B_{\text{trans}}
\end{bmatrix} \in \mathbb{R}^{6 \times 6}
\end{equation}

\textbf{解耦的虚拟刚度矩阵}：
\begin{equation}
K = \begin{bmatrix}
K_{\text{rot}} & 0 \\
0 & K_{\text{trans}}
\end{bmatrix} \in \mathbb{R}^{6 \times 6}
\end{equation}

\section{两种方法的对比}

\subsection{方法1：使用Twists和Wrenches（耦合）}

\textbf{方程形式}：
\begin{equation}
M\dot{V} + BV + K(S\theta) = F_{\text{ext}}
\end{equation}

\textbf{特点}：
\begin{itemize}
\item \textbf{耦合}：旋转和平移之间存在耦合
\item \textbf{矩阵结构}：$M, B, K$是$6 \times 6$的，通常非对角
\item \textbf{耦合项}：$M_{\omega v} \neq 0$，$M_{v\omega} \neq 0$（一般情况）
\item \textbf{物理意义}：更准确地描述旋转和平移之间的相互作用
\end{itemize}

\textbf{优点}：
\begin{itemize}
\item \textbf{更准确}：可以描述旋转和平移之间的耦合
\item \textbf{更通用}：适用于一般情况
\item \textbf{符合物理}：实际系统中旋转和平移往往是耦合的
\end{itemize}

\textbf{缺点}：
\begin{itemize}
\item \textbf{复杂}：需要设计$6 \times 6$的耦合矩阵
\item \textbf{参数多}：需要确定36个参数（对于每个矩阵）
\item \textbf{难以调参}：耦合项的影响难以直观理解
\end{itemize}

\subsection{方法2：解耦方式}

\textbf{方程形式}：
\begin{align}
M_{\text{trans}}\ddot{p} + B_{\text{trans}}\dot{p} + K_{\text{trans}}p &= f_{\text{ext,trans}} \\
M_{\text{rot}}\dot{\omega} + B_{\text{rot}}\omega + K_{\text{rot}}\theta_{\text{rot}} &= m_{\text{ext}}
\end{align}

\textbf{特点}：
\begin{itemize}
\item \textbf{解耦}：旋转和平移完全独立
\item \textbf{矩阵结构}：$M, B, K$是块对角的
\item \textbf{耦合项}：$M_{\omega v} = 0$，$M_{v\omega} = 0$
\item \textbf{物理意义}：假设旋转和平移互不影响
\end{itemize}

\textbf{优点}：
\begin{itemize}
\item \textbf{简单}：只需要设计两个独立的$3 \times 3$矩阵
\item \textbf{参数少}：每个矩阵只需要9个参数（总共18个）
\item \textbf{易于调参}：旋转和平移可以独立调整
\item \textbf{直观}：物理意义清晰
\end{itemize}

\textbf{缺点}：
\begin{itemize}
\item \textbf{近似}：忽略了旋转和平移之间的耦合
\item \textbf{不准确}：如果实际系统存在耦合，模型不准确
\item \textbf{限制}：只适用于耦合可以忽略的情况
\end{itemize}

\section{具体例子}

\subsection{例子1：耦合情况}

\textbf{场景}：机器人末端执行器有非对称质量分布

\textbf{虚拟质量矩阵}（耦合）：
\begin{equation}
M = \begin{bmatrix}
M_{\omega\omega} & M_{\omega v} \\
M_{v\omega} & M_{vv}
\end{bmatrix} = \begin{bmatrix}
I_{\text{rot}} & M_{\text{coupling}} \\
M_{\text{coupling}}^T & m_{\text{trans}}I
\end{bmatrix}
\end{equation}

其中$M_{\text{coupling}} \neq 0$表示旋转和平移的耦合。

\textbf{物理意义}：
\begin{itemize}
\item 当机器人旋转时，由于非对称质量分布，会产生平移加速度
\item 当机器人平移时，由于非对称质量分布，会产生旋转加速度
\end{itemize}

\textbf{使用耦合方式}：
\begin{itemize}
\item 可以准确描述这种耦合
\item 需要设计耦合项$M_{\text{coupling}}$
\end{itemize}

\textbf{使用解耦方式}：
\begin{itemize}
\item 忽略耦合项$M_{\text{coupling}} = 0$
\item 模型不准确，但实现简单
\end{itemize}

\subsection{例子2：解耦情况}

\textbf{场景}：机器人末端执行器质量分布对称，且旋转中心与质心重合

\textbf{虚拟质量矩阵}（解耦）：
\begin{equation}
M = \begin{bmatrix}
I_{\text{rot}} & 0 \\
0 & m_{\text{trans}}I
\end{bmatrix}
\end{equation}

其中非对角块为零。

\textbf{物理意义}：
\begin{itemize}
\item 旋转运动不影响平移
\item 平移运动不影响旋转
\end{itemize}

\textbf{使用解耦方式}：
\begin{itemize}
\item 准确描述系统
\item 实现简单
\end{itemize}

\textbf{使用耦合方式}：
\begin{itemize}
\item 也可以使用，但增加了不必要的复杂性
\end{itemize}

\section{指数坐标$S\theta$的解释}

\subsection{什么是指数坐标？}

\textbf{指数坐标}：SE(3)群元素的参数化表示。

\textbf{定义}：
\begin{equation}
X = \exp([S\theta])
\end{equation}

其中：
\begin{itemize}
\item $S \in \mathbb{R}^6$：screw轴（螺旋轴）
\item $\theta \in \mathbb{R}$：沿螺旋轴的位移
\item $[S\theta]$：twist的矩阵表示
\item $\exp$：SE(3)的指数映射
\end{itemize}

\textbf{更常用的形式}：
\begin{equation}
X = \exp([\xi])
\end{equation}

其中$\xi = S\theta \in \mathbb{R}^6$是twist。

\subsection{在阻抗控制中的使用}

\textbf{位置表示}：
\begin{itemize}
\item 在原始方程中：$x \in \mathbb{R}^n$（最小坐标）
\item 在twist/wrench方式中：$S\theta$（指数坐标）
\item 两者都表示位置/位姿
\end{itemize}

\textbf{刚度项}：
\begin{equation}
K(S\theta)
\end{equation}

\textbf{物理意义}：
\begin{itemize}
\item 虚拟弹簧的恢复力
\item 依赖于当前位姿（用指数坐标表示）
\item 如果偏离平衡位置，产生恢复力
\end{itemize}

\section{选择建议}

\subsection{何时使用耦合方式（Twist/Wrench）？}

\begin{itemize}
\item \textbf{系统存在明显耦合}：
\begin{itemize}
\item 非对称质量分布
\item 旋转中心与质心不重合
\item 复杂的几何形状
\end{itemize}

\item \textbf{需要高精度}：
\begin{itemize}
\item 对精度要求高
\item 耦合项的影响不可忽略
\end{itemize}

\item \textbf{有足够的建模能力}：
\begin{itemize}
\item 可以准确建模耦合项
\item 有足够的参数辨识能力
\end{itemize}
\end{itemize}

\subsection{何时使用解耦方式？}

\begin{itemize}
\item \textbf{系统近似解耦}：
\begin{itemize}
\item 质量分布对称
\item 旋转中心与质心重合
\item 耦合项很小，可以忽略
\end{itemize}

\item \textbf{实现简单优先}：
\begin{itemize}
\item 需要快速实现
\item 参数调优简单
\item 不需要高精度
\end{itemize}

\item \textbf{独立控制旋转和平移}：
\begin{itemize}
\item 旋转和平移有不同的控制目标
\item 需要独立调整参数
\end{itemize}
\end{itemize}

\section{导纳控制中的耦合问题}

\subsection{导纳控制也会遇到耦合问题}

\textbf{问题}：导纳控制中是否也会有耦合问题？

\textbf{答案}：\textcolor{red}{是的！}导纳控制也会遇到同样的耦合问题。

\subsection{导纳控制的方程}

\textbf{标准导纳控制方程}：
\begin{equation}
M\ddot{x}_d + B\dot{x} + Kx = f_{\text{ext}} \tag{11.64-admittance}
\end{equation}

\textbf{使用twists和wrenches的导纳控制方程}：
\begin{equation}
M\dot{V}_d + BV + K(S\theta) = F_{\text{ext}} \tag{11.64-admittance-twist}
\end{equation}

其中：
\begin{itemize}
\item $F_{\text{ext}} \in \mathbb{R}^6$：感测到的外力wrench（6维）
\item $V \in \mathbb{R}^6$：当前twist（6维）
\item $\dot{V}_d \in \mathbb{R}^6$：期望的twist加速度（6维）
\item $M, B, K \in \mathbb{R}^{6 \times 6}$：虚拟质量、阻尼、刚度矩阵
\end{itemize}

\subsection{耦合问题的表现}

\textbf{如果矩阵$M, B, K$是耦合的}（非对角）：
\begin{itemize}
\item 旋转外力矩$m_{\text{ext}}$会影响平移加速度$\dot{v}_d$
\item 平移外力$f_{\text{ext}}$会影响旋转加速度$\dot{\omega}_d$
\item 旋转和平移之间存在耦合
\end{itemize}

\textbf{具体例子}：

假设虚拟质量矩阵是耦合的：
\begin{equation}
M = \begin{bmatrix}
M_{\omega\omega} & M_{\omega v} \\
M_{v\omega} & M_{vv}
\end{bmatrix}
\end{equation}

其中$M_{\omega v} \neq 0$，$M_{v\omega} \neq 0$。

\textbf{导纳控制方程展开}：
\begin{align}
\begin{bmatrix}
M_{\omega\omega} & M_{\omega v} \\
M_{v\omega} & M_{vv}
\end{bmatrix}
\begin{bmatrix}
\dot{\omega}_d \\
\dot{v}_d
\end{bmatrix}
+ 
\begin{bmatrix}
B_{\omega\omega} & B_{\omega v} \\
B_{v\omega} & B_{vv}
\end{bmatrix}
\begin{bmatrix}
\omega \\
v
\end{bmatrix}
+ 
\begin{bmatrix}
K_{\omega\omega} & K_{\omega v} \\
K_{v\omega} & K_{vv}
\end{bmatrix}
\begin{bmatrix}
\theta_{\text{rot}} \\
\theta_{\text{trans}}
\end{bmatrix}
= 
\begin{bmatrix}
m_{\text{ext}} \\
f_{\text{ext}}
\end{bmatrix}
\end{align}

\textbf{耦合项的影响}：
\begin{itemize}
\item $M_{\omega v}\dot{v}_d$：平移加速度对旋转加速度的影响
\item $M_{v\omega}\dot{\omega}_d$：旋转加速度对平移加速度的影响
\item $B_{\omega v}v$：平移速度对旋转阻尼的影响
\item $B_{v\omega}\omega$：旋转速度对平移阻尼的影响
\item $K_{\omega v}\theta_{\text{trans}}$：平移位置对旋转刚度的影响
\item $K_{v\omega}\theta_{\text{rot}}$：旋转位置对平移刚度的影响
\end{itemize}

\subsection{解耦的导纳控制}

\textbf{解耦方式}：将旋转和平移分开处理。

\textbf{旋转部分}：
\begin{equation}
M_{\text{rot}}\dot{\omega}_d + B_{\text{rot}}\omega + K_{\text{rot}}\theta_{\text{rot}} = m_{\text{ext}}
\end{equation}

\textbf{平移部分}：
\begin{equation}
M_{\text{trans}}\dot{v}_d + B_{\text{trans}}v + K_{\text{trans}}\theta_{\text{trans}} = f_{\text{ext}}
\end{equation}

\textbf{解耦的矩阵结构}：
\begin{equation}
M = \begin{bmatrix}
M_{\text{rot}} & 0 \\
0 & M_{\text{trans}}
\end{bmatrix}, \quad
B = \begin{bmatrix}
B_{\text{rot}} & 0 \\
0 & B_{\text{trans}}
\end{bmatrix}, \quad
K = \begin{bmatrix}
K_{\text{rot}} & 0 \\
0 & K_{\text{trans}}
\end{bmatrix}
\end{equation}

\subsection{导纳控制 vs 阻抗控制的耦合}

\textbf{相同点}：
\begin{itemize}
\item 两者都会遇到旋转-平移耦合问题
\item 耦合的来源相同：虚拟质量、阻尼、刚度矩阵的非对角项
\item 解耦方法相同：使用块对角矩阵
\end{itemize}

\textbf{不同点}：
\begin{itemize}
\item \textbf{阻抗控制}：感测运动，产生力（耦合体现在力的产生上）
\item \textbf{导纳控制}：感测力，产生运动（耦合体现在运动的产生上）
\end{itemize}

\textbf{物理意义}：
\begin{itemize}
\item 在阻抗控制中：旋转运动可能产生平移力（如果$M_{\omega v} \neq 0$）
\item 在导纳控制中：旋转力矩可能产生平移加速度（如果$M_{\omega v} \neq 0$）
\item 两者是对偶的
\end{itemize}

\section{总结}

\subsection{关键区别}

\begin{center}
\begin{tabular}{|l|l|l|}
\hline
\textbf{特性} & \textbf{Twist/Wrench方式} & \textbf{解耦方式} \\
\hline
耦合性 & \textcolor{red}{耦合} & \textcolor{green}{解耦} \\
\hline
矩阵结构 & $6 \times 6$，一般非对角 & $6 \times 6$，块对角 \\
\hline
耦合项 & $M_{\omega v} \neq 0$ & $M_{\omega v} = 0$ \\
\hline
参数数量 & 36个（每个矩阵） & 18个（每个矩阵） \\
\hline
复杂度 & 高 & 低 \\
\hline
准确性 & 高（如果建模准确） & 中等（忽略耦合） \\
\hline
适用场景 & 存在明显耦合 & 耦合可忽略 \\
\hline
\end{tabular}
\end{center}

\subsection{导纳控制中的耦合}

\begin{enumerate}
\item \textbf{导纳控制也会遇到耦合问题}
\begin{itemize}
\item 如果虚拟质量矩阵$M$是耦合的，旋转力矩会影响平移加速度
\item 如果虚拟阻尼矩阵$B$是耦合的，旋转速度会影响平移阻尼
\item 如果虚拟刚度矩阵$K$是耦合的，旋转位置会影响平移刚度
\end{itemize}

\item \textbf{解耦方法相同}
\begin{itemize}
\item 使用块对角矩阵
\item 将旋转和平移分开处理
\item 独立设计旋转和平移的阻抗参数
\end{itemize}

\item \textbf{与阻抗控制对偶}
\begin{itemize}
\item 阻抗控制：运动$\to$力（耦合体现在力的产生）
\item 导纳控制：力$\to$运动（耦合体现在运动的产生）
\item 两者遇到相同的耦合问题
\end{itemize}
\end{enumerate}

\subsection{核心结论}

\begin{enumerate}
\item \textbf{Twist/Wrench方式}：
\begin{itemize}
\item \textcolor{red}{不解耦}
\item 旋转和平移之间存在耦合
\item 更准确，但更复杂
\end{itemize}

\item \textbf{解耦方式}：
\begin{itemize}
\item \textcolor{green}{解耦}
\item 旋转和平移完全独立
\item 更简单，但可能不够准确
\end{itemize}

\item \textbf{选择依据}：
\begin{itemize}
\item 如果系统存在明显耦合：使用Twist/Wrench方式
\item 如果耦合可以忽略：使用解耦方式
\item 如果实现简单优先：使用解耦方式
\end{itemize}
\end{enumerate}

\end{document}

